\hspace{3mm}

\centerline{\bf \huge 5 класс}

\begin{flushright}
        \scriptsize{ $-$ Если сдаваться, не попробовав, результат не изменится.}
\end{flushright}

\problem{1} Расул Владимирович придумал систему перевода цифр в буквы и обратно, в которой каждой из десяти цифр соответствует одна уникальная буква. Расул Владимирович решил поделиться своей системой с директором Элистинского технического лицея. Директор же решила воспользоваться интересной системой и сложила в столбик два слова ЭЛМШ с ЭТЛ, получив суммарно МШМЭЛ. Определите все варианты чисел, скрывающихся за ЭЛМШ с ЭТЛ, и докажите, что других вариантов быть не может.

\problem{2} Вчера вечером Адьян Валериевич загадал некоторое \textit{интересное} число. Число называется \textit{интересным}, если оно ровно в $2025$ раз больше одной из своих цифр. Найдите все \textit{интересные} числа, которые мог загадать Адьян Валериевич, и докажите, что других быть не может.

\problem{3}
Сугуру, Сукуна и Сатору решают олимпиаду ЭлМШ. Пока Сугуру решает $3$ задачи, Сукуна решает $4$, а пока Сукуна решает $3$ задачи, Сатору решает $5$ задач. Сколько задач решит Сугуру, если все втроём решили в общей сложности $9225$ задач?

\problem{4} 
Однажды Елена Михайловна выписала по кругу числа в следующем порядке:

$1, 2, 3, 4, 5, 6, 7, 9.$

За одно действие можно прибавить или отнять у двух рядом стоящих чисел по единице. Может ли Елена Михайловна получить следующие числа в следующем порядке:

а) $0, 0, 0, 0, 0, 0, 0, 2$? 

б) $2021, 2022, 2023, 2024, 2025, 2026, 2027, 2019$? 

\problem{5}
В турнире по настольному теннису участвовали $128$ человек. В каждом туре участники разбивались на пары и играли, после чего проигравшие выбывали из турнира (ничьих не бывает). После турнира, когда определился единоличный победитель, каждый из участников заявил, что в течение турнира обыграл не более одного правши. Причём все правши сказали правду, а все левши соврали. Сколько правшей могло участвовать в турнире?
